% 编译：xelatex literature-and-global-model.tex（需要 ctex 宏包，编译两遍以解析引用）
\documentclass[11pt,a4paper]{ctexart}
\usepackage[margin=2.5cm]{geometry}
\usepackage{amsmath,amssymb}
\usepackage{booktabs,tabularx}
\usepackage{enumitem}
\usepackage[hidelinks]{hyperref}
\usepackage{xurl}

\newcommand{\E}{\mathbb{E}}
\DeclareMathOperator{\Var}{Var}
\DeclareMathOperator{\Cov}{Cov}

\title{分层文献调研与全局误差传播模型}
\date{2026年11月12日}

\begin{document}
\maketitle

\begin{abstract}
本文按器件、电路、框架、应用四层，列出每层需事先阅读的论文、可借鉴的方法、需要注意的关键点，并给出每层的模型；最后把四层模型组合成一个全局模型。全局模型采用“双通道”结构：连续的小误差通道（近似误差、轻微时序错误）用矩传播处理，离散的罕见大误差通道（高位时序错误）借鉴可靠性分析中的屏蔽概率处理。文献信息来自检索摘要，标注“待核实”的条目需回原文确认。
\end{abstract}

\tableofcontents

\section{阅读优先级说明}
\begin{itemize}[leftmargin=*]
  \item \textbf{A（必读）}：方法将被直接采用或作为直接对比对象。
  \item \textbf{B（选读）}：提供背景、实验设置或补充方法。
\end{itemize}

\section{L1 器件层}

\subsection{阅读清单}
\begin{table}[h]
\centering
\small
\begin{tabularx}{\textwidth}{@{}clX@{}}
\toprule
级别 & 文献 & 读什么 \\
\midrule
A & Salami 等，DSN 2020 \cite{salami2020} & \textbf{在 3 块 ZCU102 上}做降压实验，与我们平台相同；降压实验流程、安全裕量、频率降低补偿 \\
A & Chamberlin 等，JLPEA 2024 \cite{chamberlin2024} & 超频脉动阵列 MAC 的时序错误特征；错误值分布“变化很小”的反直觉结论 \\
A & LEVAX，TCAD 2019/2020 \cite{levax} & 按输出位、依赖输入的时序错误率学习模型 \\
B & TEVoT，DAC 2020 \cite{tevot} & 逐输出的确定性时序违例分类 \\
B & Givaki 等，arXiv 2001.00053 \cite{givaki2020} & 降压导致的片上存储错误对 DNN 训练的影响 \\
B & Training for “Unstable” CNN Accelerator，arXiv 1812.01689 \cite{unstable2018} & 超频 CNN 加速器；把错误行为放进训练回路 \\
B & Marty、Yuki、Derrien（INRIA）\cite{marty} & 超频与轻量错误检测结合；错误率随温度、板卡、输入数据变化 \\
\bottomrule
\end{tabularx}
\caption{L1 阅读清单}
\end{table}

\subsection{借鉴的方法}
\begin{itemize}[leftmargin=*]
  \item Salami 等的降压实验流程与“保护带”概念：先找出无错误的最低电压，再在其下逐步测量。
  \item LEVAX 的思路：按输出位、以输入为条件建模出错率，而不是给整个单元一个总错误率。
  \item Chamberlin 等的超频扫描方法：单块板、固定布局、逐频率扫描。
\end{itemize}

\subsection{关键注意点}
\begin{enumerate}[leftmargin=*]
  \item 时序错误\textbf{依赖输入序列和数据}，同一频率下不同数据的错误率差别很大。
  \item Chamberlin 等发现错误输出的数值变化很小，与“随机位翻转”的直觉相反；我们的模型不能预设错误分布形状。
  \item 不同板卡之间差异明显（Salami 用了 3 块板）；我们只有两块不同型号的板，不能做统计推断。
  \item 温度必须记录并尽量固定。
\end{enumerate}

\subsection{本层模型 $\mathcal{M}_1$}
\begin{equation}
\mathcal{M}_1:\ (f, V, T, \text{输入跳变类别}) \ \mapsto\ \{P_k\}_{k},\ \Cov(\delta_j, \delta_k),
\end{equation}
以 probit/逻辑回归拟合 $P_k$ 随时钟周期裕量的变化，未测量点插值。输出接口 I1。

\section{L2 电路层}

\subsection{阅读清单}
\begin{table}[h]
\centering
\small
\begin{tabularx}{\textwidth}{@{}clX@{}}
\toprule
级别 & 文献 & 读什么 \\
\midrule
A & EvoApprox8b，DATE 2017 \cite{evoapprox} & 近似加法器、乘法器库及其误差指标；我们的单元来源 \\
A & Jiang 等，Proc.\ IEEE 2020 \cite{jiang2020} & 近似运算电路综述与统一评价；结论：复杂计算对加法误差比乘法误差更敏感 \\
A & Zervakis 等，DAC 2021 \cite{zervakis2021} & 用控制变量法抵消近似乘法的误差（含偏置），无需重训练 \\
A & Mrazek 等，ICCAD 2016 \cite{mrazek2016}（待核实） & 按数据分布加权的误差指标（WMED） \\
B & Jiang 等，ACM JETC 2017 \cite{jiang2017} & 近似运算电路的分类与比较 \\
B & 阵列乘法器误差分析，arXiv 1908.01343 \cite{arraymult} & 误差随输入变化；平均指标会掩盖大误差 \\
\bottomrule
\end{tabularx}
\caption{L2 阅读清单}
\end{table}

\subsection{借鉴的方法}
\begin{itemize}[leftmargin=*]
  \item 用 EvoApprox8b 的 C 模型穷举生成误差表 $e(a,b)$。
  \item WMED 的思想：误差指标按真实数据分布加权。我们进一步按权重取条件（$\mu(a)$、$\sigma^2(a)$）。
  \item 控制变量法：既是一种偏置补偿手段，也说明偏置可以被解析地估计和抵消，可作为我们模型的应用案例。
\end{itemize}

\subsection{关键注意点}
\begin{enumerate}[leftmargin=*]
  \item EvoApprox8b 中许多乘法器是无符号的；有符号数的处理方式决定偏置能否抵消。
  \item 近似单元的关键路径与精确单元不同，时序错误必须在近似单元本身上测。
  \item 近似误差与时序错误可能相关，叠加时要检验协方差。
\end{enumerate}

\subsection{本层模型 $\mathcal{M}_2$}
\begin{equation}
\mathcal{M}_2:\ \bigl(e_{\text{a}}(a,b),\ p_x,\ \mathcal{M}_1\bigr) \ \mapsto\ \bigl\{\mu(a),\ \sigma^2(a)\bigr\}_a,\ \ \lambda,\ \ F_{\text{big}},
\end{equation}
其中 $\mu(a)$、$\sigma^2(a)$ 描述小误差（近似误差与低位时序错误），$\lambda$ 为每次运算发生大误差（高位时序错误）的概率，$F_{\text{big}}$ 为大误差的幅度分布。输出接口 I2。

\section{L3 框架层}

\subsection{阅读清单}
\begin{table}[h]
\centering
\small
\begin{tabularx}{\textwidth}{@{}clX@{}}
\toprule
级别 & 文献 & 读什么 \\
\midrule
A & Gast \& Roth，CVPR 2018 \cite{gast2018} & 假设密度滤波（ADF）：均值、方差逐层传播的完整配方，含线性层、池化、最大值的闭式矩 \\
A & Wu 等，ICLR 2019 \cite{wu2019}（待核实） & 带协方差的确定性矩传播；ReLU 的闭式矩 \\
A & 近似乘法器误差传播，arXiv 2512.06537 \cite{amierr} & 误差矩到 GEMM 输出失真的解析模型；偏置主导的结论，需与我们的判据对照 \\
A & Lin 等，ICML 2016 \cite{lin2016}（待核实） & 量化噪声的信噪比在层间可加的分析方法 \\
A & AdaPT，TCAD 2022 \cite{adapt}；TFApprox/tf-approximate \cite{tfapprox} & 仿真引擎的实现参考 \\
B & A2Q，ICCV 2023 \cite{a2q} & 累加器位宽的溢出界（权重 L1 范数约束） \\
B & ApproxTrain \cite{approxtrain} & GPU 上的近似乘法仿真加速 \\
\bottomrule
\end{tabularx}
\caption{L3 阅读清单}
\end{table}

\subsection{借鉴的方法}
\begin{itemize}[leftmargin=*]
  \item \textbf{ADF 矩传播（Gast \& Roth）作为解析引擎的骨架}：线性层 $\mu_{\text{out}} = W\mu_{\text{in}} + b$，方差用平方权重传播；ReLU、最大池化用正态闭式矩。
  \item \textbf{协方差传播（Wu 等）}：在需要时保留协方差，而不是只传对角方差。
  \item \textbf{A2Q 的溢出界}：判断累加器位宽是否会溢出，决定是否需要建模饱和。
  \item \textbf{AdaPT/tf-approximate}：查找表替换乘法的仿真实现。
\end{itemize}

\subsection{关键注意点}
\begin{enumerate}[leftmargin=*]
  \item ADF 原本用于传播输入不确定性，误差来源只在输入端；我们需要在\textbf{每一层注入新的误差}，并处理新误差与传播误差的相关性。
  \item ADF 通常假设各元素独立（只传对角方差）；共享权重导致的输出间相关会被漏掉，需用协方差摘要补上。
  \item 原始 ADF 没有整数重量化；必须加入舍入掩蔽（见数学文档）。
  \item 推理时折叠的 BatchNorm 不吸收偏置。
\end{enumerate}

\subsection{本层模型 $\mathcal{M}_3$}
在 ADF 基础上扩展的逐层递推：
\begin{align}
m_{l+1} &= \mathcal{G}_l^{\mu}\bigl(m_l, \Sigma_l\bigr) + \mu_{\varepsilon}^{(l)}, \\
\Sigma_{l+1} &= \mathcal{G}_l^{\Sigma}\bigl(m_l, \Sigma_l\bigr) + \Sigma_{\varepsilon}^{(l)} + \text{交叉项},
\end{align}
其中 $\mathcal{G}_l$ 是第 $l$ 层的矩映射（线性层精确、ReLU 与池化用闭式矩、重量化用舍入掩蔽公式），$\Sigma$ 以“每通道方差 + 空间相关函数”的形式摘要。输出接口 I3。

\section{L4 应用层}

\subsection{阅读清单}
\begin{table}[h]
\centering
\small
\begin{tabularx}{\textwidth}{@{}clX@{}}
\toprule
级别 & 文献 & 读什么 \\
\midrule
A & Sakr、Kim、Shanbhag，ICML 2017 \cite{sakr2017} & \textbf{浮点与定点网络的失配概率}及其上界；噪声增益可由反向传播计算 \\
A & Li 等，SC 2017 \cite{li2017} & DNN 加速器中的错误传播：只有高位出错会导致输出错误；层越深越敏感；选择性加固 \\
A & ALWANN，ICCAD 2019 \cite{alwann} & 逐层选择近似乘法器、不重训练的设计探索 \\
A & HAWQ-V2/V3 \cite{hawq} & 基于 Hessian 迹的层敏感度；用整数线性规划选择每层配置 \\
\bottomrule
\end{tabularx}
\caption{L4 阅读清单}
\end{table}

\subsection{借鉴的方法}
\begin{itemize}[leftmargin=*]
  \item \textbf{Sakr 等的失配概率}：我们在数学文档中提出的“基于间隔的翻转概率”与它是同一思路，应以 Sakr 等为基础并引用；其“噪声增益”即数学文档中设计探索公式里的 $g_l$。
  \item \textbf{Li 等的位敏感度结论}：只有高位错误导致输出错误，支持全局模型中把大误差单独作为一个通道处理。
  \item \textbf{HAWQ-V3 的整数线性规划}：把逐层配置选择写成整数规划，作为我们设计探索的求解方法。
  \item \textbf{ALWANN}：作为设计探索的直接对比基线。
\end{itemize}

\subsection{关键注意点}
\begin{enumerate}[leftmargin=*]
  \item Sakr 等分析的是量化噪声（零均值、近似均匀），近似误差有偏置；需要把偏置项加入失配概率。
  \item 共模偏置在 softmax/argmax 中被抵消，只有 logits 之间误差的差起作用。
  \item Hessian 敏感度是对损失函数的二阶近似，与精度下降并不等价，需实测验证。
  \item 精度测量本身有约 0.3\% 的标准误差，应使用配对比较。
\end{enumerate}

\subsection{本层模型 $\mathcal{M}_4$}
\begin{equation}
\mathcal{M}_4:\ \text{I3} \ \mapsto\ P_{\text{flip},s} \ \mapsto\ \widehat{\Delta\text{Acc}},
\end{equation}
在 Sakr 等失配概率的基础上加入偏置项与 logits 间协方差，并用缩放因子 $\kappa$ 标定尾部。

\section{全局模型}

\subsection{动机：为什么需要两个通道}
\begin{itemize}[leftmargin=*]
  \item 近似误差每次运算都有、幅度小，适合用均值与方差描述，并经矩传播计算（Gast \& Roth、Sakr 等的思路）。
  \item 高位时序错误罕见但幅度大，一旦发生往往直接改变输出（Li 等的结论）；均值与方差描述不了，更适合用“发生概率 $\times$ 被屏蔽的概率”描述（Vallero 等 \cite{vallero2016} 的思路）。
\end{itemize}
因此全局模型由\textbf{小误差通道}与\textbf{大误差通道}组成，两者最后在应用层合并。

\subsection{输入}
系统配置
\begin{equation}
c = \bigl(f,\ V,\ T,\ \{m_l\}_{l=1}^{L},\ \text{架构参数}\bigr),
\end{equation}
以及模型、测试数据、各层数据分布。

\subsection{计算步骤}

\paragraph{第 1 步：器件层.}
由 $\mathcal{M}_1$ 得到每个输出位的出错概率 $\{P_k(f,V,T)\}$。

\paragraph{第 2 步：电路层（拆分为两个通道）.}
按位权把时序错误分为低位与高位两部分：
\begin{itemize}[leftmargin=*]
  \item 低位时序错误与近似误差合并，得到小误差通道的条件矩 $\mu_l(a)$、$\sigma_l^2(a)$；
  \item 高位时序错误给出每次运算发生大误差的概率 $\lambda_l$ 和幅度分布 $F_{\text{big},l}$。
\end{itemize}
低位与高位的分界位 $k^*$ 是模型参数，按“该位出错能否被小误差通道的正态近似覆盖”选取，并通过实验确定。

\paragraph{第 3 步：框架层.}
\begin{itemize}[leftmargin=*]
  \item \textbf{小误差通道}：用扩展 ADF（$\mathcal{M}_3$）逐层传播均值与协方差摘要，得到输出层 logits 误差的统计量。
  \item \textbf{大误差通道}：对每个网络层 $l$ 估计屏蔽概率的补，即“一次大误差导致输出改变的概率” $s_l$。$s_l$ 只与模型和数据有关，与配置无关，对每个模型用少量故障注入估计一次（对应 Vallero 等的 CPT 表征），之后反复使用。
\end{itemize}

\paragraph{第 4 步：应用层合并.}
对样本 $s$，小误差通道给出翻转概率 $P_{\text{small},s}$（$\mathcal{M}_4$）。第 $l$ 层有 $n_l$ 次运算，大误差通道导致翻转的概率为
\begin{equation}
P_{\text{big}} = 1 - \prod_{l} \bigl(1 - \lambda_l s_l\bigr)^{n_l} \approx \sum_l n_l \lambda_l s_l \quad(\text{当乘积很小时}).
\end{equation}
假设两个通道独立：
\begin{equation}
P_{\text{flip},s} = 1 - \bigl(1 - P_{\text{small},s}\bigr)\bigl(1 - P_{\text{big}}\bigr),
\end{equation}
\begin{equation}
\widehat{\Delta\text{Acc}} = \frac{1}{S}\sum_{s \in \mathcal{C}} P_{\text{flip},s} \ -\ (\text{原本错误、被改对的小项}).
\end{equation}

\subsection{需要标定的参数}
\begin{table}[h]
\centering
\small
\begin{tabularx}{\textwidth}{@{}lXl@{}}
\toprule
参数 & 含义 & 来源 \\
\midrule
$\{P_k\}$ 的拟合参数 & 时序错误模型 & L1 实测 \\
$k^*$ & 大小误差的分界位 & 实验确定 \\
$s_l$ & 每层一次大误差导致输出改变的概率 & 每个模型一次故障注入 \\
$\kappa$ & 小误差通道尾部的缩放因子 & 少量实测配置 \\
\bottomrule
\end{tabularx}
\caption{全局模型的标定参数}
\end{table}

\subsection{计算代价}
\begin{itemize}[leftmargin=*]
  \item 小误差通道：一次带矩的前向传播，约为普通前向传播的 2--3 倍。
  \item 大误差通道：$s_l$ 预先算好，预测时只需求和。
  \item 因此单个配置的预测时间在秒级，满足验收标准中“$\le 1$ 秒、比逐比特仿真快 100 倍以上”的目标。
\end{itemize}

\subsection{与已有工作的关系}
\begin{table}[h]
\centering
\small
\begin{tabularx}{\textwidth}{@{}lX@{}}
\toprule
已有工作 & 本模型的扩展 \\
\midrule
Vallero 等 \cite{vallero2016} & 只有离散故障通道；本模型加入连续的小误差通道，并以数据分布为条件 \\
Gast \& Roth \cite{gast2018} & 只在输入端引入不确定性；本模型在每层注入误差，加入重量化与相关性 \\
Sakr 等 \cite{sakr2017} & 只考虑零均值量化噪声；本模型加入偏置、协方差与大误差通道 \\
2512.06537 \cite{amierr} & 只到 GEMM 输出失真；本模型逐层传播到最终精度，并包含器件层时序错误 \\
Li 等 \cite{li2017} & 用故障注入统计敏感度；本模型把它作为大误差通道的参数 $s_l$，与解析的小误差通道结合 \\
\bottomrule
\end{tabularx}
\caption{全局模型与已有工作的关系}
\end{table}

\subsection{模型的主要假设（需逐一检验）}
\begin{enumerate}[leftmargin=*]
  \item 两个通道相互独立。
  \item 小误差通道的正态近似（经 $\kappa$ 标定后）在应用层足够准确。
  \item 大误差之间不发生叠加（大误差足够罕见）。
  \item $s_l$ 与配置无关，只依赖模型和数据。
  \item 新注入误差与传播误差之间的交叉项可以忽略，或可由摘要统计量近似。
\end{enumerate}

\subsection{验证方案}
按实验方案阶段 5：在至少 20 个未参与标定的配置上，比较全局模型预测与板上实测；并做消融实验——去掉大误差通道、去掉协方差、改用均匀分布误差矩、去掉舍入掩蔽——量化每个组成部分对预测精度的贡献。

\section{阅读顺序建议}
\begin{enumerate}[leftmargin=*]
  \item 第 1 周：Vallero 等（已读）、Sakr 等、Gast \& Roth——全局模型的三根支柱。
  \item 第 2 周：Li 等、2512.06537、Salami 等——大误差通道、对比基线、平台先例。
  \item 第 3 周：EvoApprox8b、Jiang 等综述、Zervakis 等、Chamberlin 等——电路层与器件层细节。
  \item 之后按需：ALWANN、HAWQ、A2Q、LEVAX、AdaPT 等。
\end{enumerate}

\begin{thebibliography}{99}
\bibitem{salami2020} B.~Salami et al. An experimental study of reduced-voltage operation in modern FPGAs for neural network acceleration. DSN, 2020. arXiv:2005.03451.
\bibitem{chamberlin2024} Chamberlin et al. Understanding timing error characteristics from overclocked systolic multiply--accumulate arrays in FPGAs. Journal of Low Power Electronics and Applications, 2024. \url{https://doaj.org/article/459c4dfd9fb149b5ad98690e9a77cb02}
\bibitem{levax} LEVAX: input-aware learning-based error model of voltage-scaled functional units. IEEE TCAD, 2019/2020.
\bibitem{tevot} TEVoT: timing error modeling of voltage-scaled circuits. DAC, 2020.
\bibitem{givaki2020} K.~Givaki et al. On the resilience of deep learning for reduced-voltage FPGAs. arXiv:2001.00053.
\bibitem{unstable2018} Training for ``unstable'' CNN accelerator: a case study on FPGA. arXiv:1812.01689.
\bibitem{marty} T.~Marty, T.~Yuki, S.~Derrien. Algorithm-level timing speculation for fault-tolerant CNN accelerators. INRIA. \url{https://core.ac.uk/download/163070624.pdf}
\bibitem{evoapprox} V.~Mrazek et al. EvoApprox8b: library of approximate adders and multipliers for circuit design and benchmarking of approximation methods. DATE, 2017.
\bibitem{jiang2020} H.~Jiang, F.\,J.\,H. Santiago, H.~Mo, L.~Liu, J.~Han. Approximate arithmetic circuits: a survey, characterization, and recent applications. Proceedings of the IEEE, 108(12):2108--2135, 2020.
\bibitem{zervakis2021} G.~Zervakis et al. Control variate approximation for DNN accelerators. DAC, 2021. arXiv:2102.09642.
\bibitem{mrazek2016} V.~Mrazek et al. Design of power-efficient approximate multipliers for approximate artificial neural networks. ICCAD, 2016.（待核实）
\bibitem{jiang2017} H.~Jiang, C.~Liu, L.~Liu, F.~Lombardi, J.~Han. A review, classification, and comparative evaluation of approximate arithmetic circuits. ACM JETC, 2017. DOI: 10.1145/3094124.
\bibitem{arraymult} Error analysis of approximate array multipliers. arXiv:1908.01343.
\bibitem{gast2018} J.~Gast, S.~Roth. Lightweight probabilistic deep networks. CVPR, 2018. arXiv:1805.11327.
\bibitem{wu2019} A.~Wu et al. Deterministic variational inference for robust Bayesian neural networks. ICLR, 2019.（待核实）
\bibitem{amierr} Approximate multiplier induced error propagation in deep neural networks. arXiv:2512.06537, 2025.
\bibitem{lin2016} D.~Lin, S.~Talathi, V.\,S.~Annapureddy. Fixed point quantization of deep convolutional networks. ICML, 2016.（待核实）
\bibitem{adapt} AdaPT: fast emulation of approximate DNN accelerators in PyTorch. IEEE TCAD, 2022. arXiv:2203.04071.
\bibitem{tfapprox} tf-approximate. \url{https://github.com/ehw-fit/tf-approximate}
\bibitem{a2q} I.~Colbert, A.~Pappalardo, J.~Petri-Koenig. A2Q: accumulator-aware quantization with guaranteed overflow avoidance. ICCV, 2023. arXiv:2308.13504.
\bibitem{approxtrain} ApproxTrain: fast simulation of approximate multipliers for DNN training and inference. arXiv:2209.04161.
\bibitem{sakr2017} C.~Sakr, Y.~Kim, N.~Shanbhag. Analytical guarantees on numerical precision of deep neural networks. ICML, 2017. PMLR 70:3007--3016.
\bibitem{li2017} G.~Li, S.\,K.\,S.~Hari, M.~Sullivan, T.~Tsai, K.~Pattabiraman, J.~Emer, S.\,W.~Keckler. Understanding error propagation in deep learning neural network (DNN) accelerators and applications. SC, 2017. DOI: 10.1145/3126908.3126964.
\bibitem{alwann} V.~Mrazek, Z.~Vasicek, L.~Sekanina, M.\,A.~Hanif, M.~Shafique. ALWANN: automatic layer-wise approximation of deep neural network accelerators without retraining. ICCAD, 2019. arXiv:1907.07229.
\bibitem{hawq} Z.~Yao et al. HAWQ-V3: dyadic neural network quantization. ICML, 2021. arXiv:2011.10680.
\bibitem{vallero2016} A.~Vallero et al. Cross-layer system reliability assessment framework for hardware faults. ITC, 2016. DOI: 10.1109/TEST.2016.7805863.
\end{thebibliography}

\end{document}
